%==============================================================================
%  CHAPTER 3 — SYSTEM DESIGN
%==============================================================================
\chapter{System Design}
\label{chap:requirements}
\label{chap:design}

\section{Introduction}
This chapter designs the system as a whole, before the next chapter details
its realization sprint by sprint. It presents the functional view (actors,
requirements, and use cases), the global architecture in its logical and
physical forms, the organization of the two pipelines, the detailed dynamic
design through sequence diagrams, and the data model.

% ---------------------------------------------------------------------------
\section{Functional View}
\label{sec:functional-view}

\subsection{Actors}
Four actors interact with the system:
\begin{itemize}[leftmargin=1.4em]
  \item \textbf{Platform user (data scientist)} --- the primary actor: manages
        projects, uploads code and data, triggers and follows training runs,
        compares experiments, registers, promotes, deploys, and tests models,
        and manages the \gls{api} keys that expose a deployment publicly. No
        infrastructure knowledge is assumed: every cluster interaction happens
        behind a button.
  \item \textbf{External API client} --- a third-party website, mobile
        application, or backend that calls a deployed model's public endpoint
        with an \gls{api} key. It has no platform account; its whole contract
        is one authenticated \gls{http} call.
  \item \textbf{Administrator} --- inherits all user capabilities and, in
        addition, manages accounts and roles and oversees every user's
        projects, pipelines, and deployments.
  \item \textbf{System (automated actors)} --- the training pipeline and the
        model server act on the user's behalf once triggered; they appear as
        participants in the sequence diagrams.
\end{itemize}

\subsection{Functional Requirements}
The platform must provide the following functional capabilities, grouped by
module:
\begin{itemize}[leftmargin=1.4em]
  \item \textbf{Authentication and sessions}: sign up, sign in, refresh the
        access token, and retrieve the current user's profile and role.
  \item \textbf{Project management}: create, list, update, and delete projects;
        bind each project to its own tracking experiment; enforce strict
        per-owner isolation (foreign resources return \texttt{404}).
  \item \textbf{Upload and code handling}: upload code
        (\texttt{.py}, \texttt{.ipynb}, \texttt{.zip}) and datasets;
        auto-detect the entry script; convert notebooks to runnable scripts;
        run a pre-flight static analysis that warns about common pitfalls.
  \item \textbf{Training execution}: trigger a training run in an isolated
        cluster pod, stream its logs live to the browser, and install missing
        dependencies automatically.
  \item \textbf{Experiment tracking}: capture parameters, metrics, and
        artifacts of every run automatically; browse and compare runs; open a
        run-details view with an artifact browser (list, preview, download).
  \item \textbf{Model registry}: register a run as a versioned model, list its
        versions, and promote them across Staging / Production / Archived.
  \item \textbf{Deployment}: deploy a registered version as a live service in
        one click, detect the model framework automatically
        (scikit-learn / XGBoost), track the deployment status and endpoint, and
        delete a deployment.
  \item \textbf{Prediction}: test predictions from the interface, with readable
        error reporting when an input or the model server fails.
  \item \textbf{Public API and keys}: create, list, and revoke per-deployment
        \gls{api} keys (stored hashed); serve a public, key-authenticated
        prediction endpoint; generate ready-to-paste code snippets.
  \item \textbf{Monitoring}: display live cluster and per-pod resource metrics,
        and a per-project serving view with request volume, error rate, latency,
        and timeline, traffic split by origin (tester versus public \gls{api}),
        and a per-deployment breakdown including each model's live \gls{cpu} and
        memory usage.
  \item \textbf{Insights}: track each model's evolution across versions with
        evaluation \glspl{kpi}, a health verdict, and version/stage progression.
  \item \textbf{Diagnostics}: capture and surface the failure reason of a
        failed run so the user understands \emph{why} it failed.
  \item \textbf{Administration}: manage user accounts (activate, deactivate,
        promote, demote, purge, delete), consult platform-wide statistics, and
        oversee every user's projects, pipelines, and deployments
        (administrator only).
\end{itemize}

\subsection{Non-Functional Requirements}
\begin{itemize}[leftmargin=1.4em]
  \item \textbf{Security}: stateless \gls{jwt} authentication~\cite{rfc7519};
        passwords stored as bcrypt hashes and \gls{api} keys as \gls{sha}-256
        hashes~\cite{owasp_api}; strict ownership checks; administrator-only
        endpoints.
  \item \textbf{Performance}: a fully asynchronous backend, so slow external
        calls never block other requests.
  \item \textbf{Scalability and isolation}: every training job and model
        server runs in its own resource-limited pod.
  \item \textbf{Portability}: identical behavior on a local KinD cluster and
        on \gls{aks}, everything expressed as Kubernetes resources.
  \item \textbf{Reproducibility}: code, parameters, metrics, and artifacts of
        every run captured automatically.
  \item \textbf{Usability}: the whole lifecycle from one web interface; no
        kubectl, no \gls{yaml}, no cloud console.
\end{itemize}

\subsection{Global Use-Case Diagram}
\label{sec:usecases}
Figure~\ref{fig:usecase-global} summarizes the interactions of the three human
actors with the system.
\begin{figure}[H]\centering
  \includegraphics[width=0.95\textwidth,height=0.48\textheight,keepaspectratio]{usecase_global}
  \caption{Global use-case diagram of the platform}
  \label{fig:usecase-global}
\end{figure}

\subsection{Use-Case Refinements}
Tables~\ref{tab:uc-signin} to~\ref{tab:uc-public} refine one central use case
per lifecycle area; the remaining flows are detailed by the sequence diagrams
of Section~\ref{sec:seq-diagrams}.

\begin{table}[H]\centering
\renewcommand{\arraystretch}{1.2}
\caption{Textual description of the ``Sign in'' use case}\label{tab:uc-signin}
\begin{tabularx}{\textwidth}{|>{\bfseries}p{3.2cm}|X|}
\hline
Use case & Sign in\\ \hline
Actor & Platform user\\ \hline
Precondition & The user already has a registered account.\\ \hline
Post-condition & The user is authenticated; an access/refresh token pair is issued.\\ \hline
Nominal scenario & 1.~The user submits email and password.\newline 2.~The system verifies the password against its stored hash.\newline 3.~The system issues a \gls{jwt} access token and a refresh token.\newline 4.~The user is redirected to the dashboard.\\ \hline
Exceptions & Invalid credentials: the request is rejected with a clear error.\\ \hline
\end{tabularx}
\end{table}

\begin{table}[H]\centering
\renewcommand{\arraystretch}{1.2}
\caption{Textual description of the ``Trigger a training run'' use case}\label{tab:uc-train}
\begin{tabularx}{\textwidth}{|>{\bfseries}p{3.2cm}|X|}
\hline
Use case & Trigger a training run\\ \hline
Actor & Platform user\\ \hline
Precondition & The user is authenticated and has uploaded code (and optionally a dataset) to a project.\\ \hline
Post-condition & A pipeline run executes on the cluster; its metrics, parameters, and artifacts are recorded.\\ \hline
Nominal scenario & 1.~The user clicks ``Run code''.\newline 2.~The backend submits a pipeline to the orchestrator.\newline 3.~An isolated pod downloads the code, installs dependencies, and runs the entry script.\newline 4.~The autolog wrapper records results to the tracking server.\newline 5.~The user follows the live logs and the captured metrics.\\ \hline
Exceptions & The script fails: the error is captured and surfaced to the user; no model is registered.\\ \hline
\end{tabularx}
\end{table}

\begin{table}[H]\centering
\renewcommand{\arraystretch}{1.2}
\caption{Textual description of the ``Deploy a model'' use case}\label{tab:uc-deploy}
\begin{tabularx}{\textwidth}{|>{\bfseries}p{3.2cm}|X|}
\hline
Use case & Deploy a model\\ \hline
Actor & Platform user\\ \hline
Precondition & A model version is registered in the registry.\\ \hline
Post-condition & A live prediction service is running and reachable on the cluster.\\ \hline
Nominal scenario & 1.~The user clicks ``Deploy'' on a registered model.\newline 2.~The backend detects the model's framework and creates an \texttt{InferenceService}.\newline 3.~The serving operator reconciles it into a Deployment, Service, and Ingress.\newline 4.~The model server loads the artifact from object storage.\newline 5.~The status and endpoint are reported back to the user.\\ \hline
Exceptions & The service fails to become ready: the deployment is marked failed and the user is notified.\\ \hline
\end{tabularx}
\end{table}

\begin{table}[H]\centering
\renewcommand{\arraystretch}{1.2}
\caption{Textual description of the ``Call the public prediction endpoint'' use case}\label{tab:uc-public}
\begin{tabularx}{\textwidth}{|>{\bfseries}p{3.2cm}|X|}
\hline
Use case & Call the public prediction endpoint\\ \hline
Actor & External \gls{api} client\\ \hline
Precondition & The model is deployed and the client holds a valid \gls{api} key.\\ \hline
Post-condition & The client receives a prediction from the deployed model.\\ \hline
Nominal scenario & 1.~The client sends a request with an \texttt{Authorization: Bearer} key.\newline 2.~The backend hashes the key and looks it up, scoped to the deployment.\newline 3.~The backend confirms the deployment is ready.\newline 4.~The request is forwarded to the model server.\newline 5.~The prediction is returned to the client.\\ \hline
Exceptions & Invalid or revoked key: the backend returns \texttt{401}.\\ \hline
\end{tabularx}
\end{table}

% ---------------------------------------------------------------------------
\section{Global Architecture}
\label{sec:architecture}

\subsection{Logical Architecture}
\label{sec:app-architecture}
The application adopts a \textbf{four-tier layered architecture}
(Figure~\ref{fig:layered-arch}):
\begin{itemize}[leftmargin=1.4em]
  \item \textbf{Presentation}: an Angular single-page application consuming
        the backend's \gls{rest} \gls{api};
  \item \textbf{API}: FastAPI routers that validate requests and enforce
        authentication and ownership;
  \item \textbf{Service}: the business logic orchestrating the cluster systems
        (pipelines, MLflow, KServe, object storage);
  \item \textbf{Data}: SQLAlchemy models mapped onto PostgreSQL.
\end{itemize}

\begin{figure}[H]\centering
  \includegraphics[width=\textwidth,keepaspectratio]{logical_arch}
  \caption{Four-tier layered architecture and the cluster systems it orchestrates}
  \label{fig:layered-arch}
\end{figure}

The layered style was preferred over microservices for proportionality: the
platform already delegates its heavyweight concerns (pipeline execution, model
serving, artifact storage) to dedicated cluster components. The result is one
backend image and one frontend image, each independently deployable.

\subsection{Physical Architecture}
\label{sec:physical}
Everything runs on one Kubernetes cluster --- \gls{aks} in production, KinD
locally, same code. Public traffic follows the path \emph{browser
$\rightarrow$ Cloudflare edge (\gls{tls}) $\rightarrow$ tunnel agent
$\rightarrow$ frontend $\rightarrow$ backend $\rightarrow$ cluster systems}.
Figure~\ref{fig:physical-arch} shows the topology: the platform's own
workloads run as long-lived deployments, while training jobs and model servers
are created and destroyed on demand, each in its own resource-limited pod.

\begin{figure}[H]\centering
  \includegraphics[width=\textwidth,height=0.55\textheight,keepaspectratio]{physical_arch}
  \caption{Physical architecture of the platform on the Kubernetes cluster}
  \label{fig:physical-arch}
\end{figure}

% ---------------------------------------------------------------------------
\section{Pipeline Organization}
\label{sec:pipeline-org}
Two flows structure the whole system: a training path and a serving path. They
meet only in the object store and the registry --- training produces artifacts
and versions, serving consumes them --- so each half can evolve, scale, and
fail independently.

The \textbf{training path} starts from an upload: the backend stores code and
dataset in the object store, then submits a \gls{kfp} pipeline whose single
step runs in an isolated pod and executes the user's code under an autolog
wrapper (Figure~\ref{fig:pipeline-flow}).

\begin{figure}[H]\centering
\resizebox{\textwidth}{!}{%
\begin{tikzpicture}
  \node[stagenode,fill=diagblue!12] (trig) {Trigger\\run};
  \node[stepnode,right=1.5cm of trig] (dl) {Download\\code + data};
  \node[stepnode,right=0.5cm of dl] (dep) {Install\\dependencies};
  \node[stepnode,right=0.5cm of dep] (conv) {Convert\\notebook};
  \node[stepnode,right=0.5cm of conv] (run) {Autolog\\runner + run};
  \node[storenode,right=1.5cm of run] (ml) {
\includegraphics[height=0.5cm]{mlflow}\\metrics \textperiodcentered{} params\\artifacts};
  \node[stagenode,below=1.1cm of conv,fill=diaggreen!10,draw=diaggreen!60] (logs) {Live logs\\$\rightarrow$ UI};
  \begin{scope}[on background layer]
    \node[clusterbox,fit=(dl)(dep)(conv)(run),
          label={[font=\scriptsize\bfseries,text=diagblue]above:{Isolated pod (Kubernetes)}}] (pod) {};
  \end{scope}
  \draw[flowarrow] (trig) -- (dl);
  \draw[flowarrow] (dl) -- (dep);
  \draw[flowarrow] (dep) -- (conv);
  \draw[flowarrow] (conv) -- (run);
  \draw[flowarrow] (run) -- (ml);
  \draw[flowarrow,draw=diaggreen!70] (pod.south) -- (logs);
\end{tikzpicture}}
\caption{Training path: from trigger to captured experiment}
\label{fig:pipeline-flow}
\end{figure}

The \textbf{serving path} starts from the registry: a registered model version
is deployed by creating one declarative \texttt{InferenceService} resource,
which the serving operator reconciles into a running model server loading its
artifact from the object store (Figure~\ref{fig:deploy-flow}).

\begin{figure}[H]\centering
\resizebox{\textwidth}{!}{%
\begin{tikzpicture}
  \node[stagenode,fill=diagblue!12] (reg) {Register\\model version};
  \node[stepnode,right=0.7cm of reg] (isvc) {Backend creates\\InferenceService};
  \node[stepnode,right=0.7cm of isvc] (op) {
\includegraphics[height=0.45cm]{kserve}\\KServe operator\\reconciles};
  \node[stepnode,right=0.7cm of op,minimum height=2.4cm] (res) {Deployment\\Service\\Ingress};
  \node[storenode,right=0.7cm of res] (srv) {Model server\\loads artifact\\
\includegraphics[height=0.4cm]{minio}};
  \node[stagenode,right=0.7cm of srv,fill=diaggreen!12,draw=diaggreen!60] (ready) {READY\\$\rightarrow$ Predict};
  \draw[flowarrow] (reg) -- (isvc);
  \draw[flowarrow] (isvc) -- (op);
  \draw[flowarrow] (op) -- (res);
  \draw[flowarrow] (res) -- (srv);
  \draw[flowarrow] (srv) -- (ready);
\end{tikzpicture}}
\caption{Serving path: from a registered model to a live prediction service}
\label{fig:deploy-flow}
\end{figure}

% ---------------------------------------------------------------------------
\section{Sequence Diagrams}
\label{sec:seq-diagrams}
The dynamic behavior of the main flows is detailed by sequence diagrams,
grouped by lifecycle area.

\subsection{Authentication}
Figures~\ref{fig:seq-signup} and~\ref{fig:seq-login} detail the sign-up and
login flows on which every later interaction depends.
\begin{figure}[H]\centering
  \includegraphics[width=0.92\textwidth,height=0.42\textheight,keepaspectratio]{seq_signup}
  \caption{Sequence diagram of the sign-up flow}
  \label{fig:seq-signup}
\end{figure}
\begin{figure}[H]\centering
  \includegraphics[width=0.92\textwidth,height=0.42\textheight,keepaspectratio]{seq_login}
  \caption{Sequence diagram of the login flow}
  \label{fig:seq-login}
\end{figure}

\subsection{Upload and Training}
Figure~\ref{fig:seq-upload} details the upload flow that feeds the pipeline,
and Figure~\ref{fig:seq-training} the training-run flow itself.
\begin{figure}[H]\centering
  \includegraphics[width=0.92\textwidth,height=0.42\textheight,keepaspectratio]{seq_upload}
  \caption{Sequence diagram of uploading code and data to a project}
  \label{fig:seq-upload}
\end{figure}
\begin{figure}[H]\centering
  \includegraphics[width=0.92\textwidth,height=0.42\textheight,keepaspectratio]{seq_training}
  \caption{Sequence diagram of triggering a custom-code training run}
  \label{fig:seq-training}
\end{figure}

\subsection{Registry and Deployment}
Figure~\ref{fig:seq-promote} details the promotion flow, and
Figure~\ref{fig:seq-deploy} the deploy-and-predict flow.
\begin{figure}[H]\centering
  \includegraphics[width=0.92\textwidth,height=0.42\textheight,keepaspectratio]{seq_promote}
  \caption{Sequence diagram of promoting a model version}
  \label{fig:seq-promote}
\end{figure}
\begin{figure}[H]\centering
  \includegraphics[width=0.92\textwidth,height=0.42\textheight,keepaspectratio]{seq_deploy}
  \caption{Sequence diagram of deploying a model and requesting a prediction}
  \label{fig:seq-deploy}
\end{figure}

\subsection{Public API and Administration}
Figure~\ref{fig:seq-public} details how an external client calls the public
endpoint, and Figure~\ref{fig:seq-admin} how an administrator manages user
accounts.
\begin{figure}[H]\centering
  \includegraphics[width=0.92\textwidth,height=0.42\textheight,keepaspectratio]{seq_public_predict}
  \caption{Sequence diagram of an external client calling the public endpoint}
  \label{fig:seq-public}
\end{figure}
\begin{figure}[H]\centering
  \includegraphics[width=0.92\textwidth,height=0.42\textheight,keepaspectratio]{seq_admin_users}
  \caption{Sequence diagram of administrator user management}
  \label{fig:seq-admin}
\end{figure}

% ---------------------------------------------------------------------------
\section{Data Model}
\label{sec:data-model}
The domain model is deliberately small. \texttt{User} owns \texttt{Project}; a
project accumulates \texttt{PipelineRun}s and \texttt{MLModel}s; a model can
become a \texttt{Deployment}, which owns its \texttt{DeploymentApiKey}s and
accumulates one \texttt{PredictionLog} row per served prediction. The
scientific content of a run (metrics, parameters, artifacts) deliberately
lives in the tracking server; the platform's database holds only identity,
ownership, lifecycle state, and telemetry. Figure~\ref{fig:class-global} shows
the global class diagram and Table~\ref{tab:entities} the entities with their
key attributes.

\begin{figure}[H]\centering
  \includegraphics[width=0.95\textwidth,height=0.55\textheight,keepaspectratio]{class_domain}
  \caption{Global class diagram of the platform}
  \label{fig:class-global}
\end{figure}

\begin{table}[H]\centering
\renewcommand{\arraystretch}{1.15}
\caption{Domain entities and key attributes}\label{tab:entities}
\small
\begin{tabularx}{\textwidth}{|l|X|}
\hline
\rowcolor{primary!15}
\textbf{Entity} & \textbf{Key attributes}\\
\hline
\texttt{User} & id; unique email; bcrypt \texttt{hashed\_password}; \texttt{role} (\texttt{user}/\texttt{admin}); \texttt{is\_active}\\ \hline
\texttt{Project} & id; name, description; \texttt{owner\_id}; \texttt{experiment\_id} of its tracking experiment\\ \hline
\texttt{PipelineRun} & id; \texttt{project\_id}; \texttt{kfp\_run\_id}; status (pending/running/succeeded/failed); timestamps\\ \hline
\texttt{MLModel} & id; \texttt{project\_id}; registered name + version; stage (Staging/Production/Archived)\\ \hline
\texttt{Deployment} & id; \texttt{model\_id}; \texttt{inference\_service\_name}; status; prediction endpoint \gls{url}\\ \hline
\texttt{DeploymentApiKey} & id; \texttt{deployment\_id}; \gls{sha}-256 \texttt{key\_hash} (plaintext never stored); display prefix; \texttt{revoked\_at}\\ \hline
\texttt{PredictionLog} & id; \texttt{deployment\_id}, \texttt{project\_id}; \texttt{latency\_ms}; \texttt{status\_code}; batch size; source (app/public); indexed timestamp\\ \hline
\end{tabularx}
\end{table}

% ---------------------------------------------------------------------------
\section{Conclusion}
This chapter fixed the design of the system: four actors and their use cases,
a layered application in front of specialized cluster systems, one training
path and one serving path meeting in the object store and the registry, the
dynamic flows as sequence diagrams, and a compact domain model. The next
chapter details how this design was realized, sprint by sprint.
